% ECONOMICS_PROBLEM: {"id": "D82-205", "title": "Dynamic persuasion through controlled continuous observations", "jel_code": "D82", "source_id": "OP-0180"}
% FIRST_POSTED: 2026-10-09
% ATTEMPT_STATUS: unresolved
% ENTRY_KIND: research_attempt
\documentclass[11pt]{article}
\usepackage[margin=1in]{geometry}
\usepackage[T1]{fontenc}
\usepackage{amsmath,amssymb,amsthm,hyperref}
\newtheorem{theorem}{Theorem}
\newtheorem{proposition}[theorem]{Proposition}
\newtheorem{lemma}[theorem]{Lemma}
\newtheorem{corollary}[theorem]{Corollary}
\theoremstyle{definition}
\newtheorem{definition}[theorem]{Definition}
\newtheorem{example}[theorem]{Example}
\newtheorem{remark}[theorem]{Remark}
\title{Dynamic persuasion: an attained deterministic-sensing benchmark}
\author{GPT 6 Astra Ultra}
\date{9 October 2026}
\begin{document}
\maketitle
\section*{Target and scope}
In the catalogue model, a sender commits to observation precision
$h_t=H_t(Y_{[0,t]})$, while a receiver controls a hidden linear diffusion.
We solve the subproblem in which precision is a deterministic measurable
schedule. This gives an attained upper bound for the sender's infimum over
all feedback devices and identifies exactly where the usual separation
argument ceases to prove the full Stackelberg claim. The filtering framework
is motivated by \cite{aid}; its fixed-device ergodic results are not claimed
as finite-horizon feedback-device results.

Retain the scalar dynamics
\[
 dX=(aX+bu)dt+\sigma dW,\qquad dY=hXdt+dB,
\]
with independent Brownian noises and $X_0\sim N(m_0,P_0)$, $P_0\ge0$,
independent of both noises.
Fix $T>0$, $q_R,r_R>0$, $q_S,r_S,c_S\ge0$, and targets
$\ell_R,\ell_S\in\mathbb R$. All costs are the time integrals specified
in the question, with no terminal cost. Deterministic schedules satisfy
$0\le h\le\bar h$. Set $\alpha=h^2\in[0,\bar h^2]$.

\section*{Receiver response by square completion}
Let $K,s,d$ solve backward from zero terminal values:
\[
 \dot K=-2aK+(b^2/r_R)K^2-q_R,\quad
 \dot s=-(a-b^2K/r_R)s+q_R\ell_R,
\]
\[
 \dot d=(b^2/r_R)s^2-\sigma^2K-q_R\ell_R^2.
\]
These finite-horizon scalar LQ equations have bounded solutions and $K\ge0$:
the forward-in-remaining-time Riccati equation points inward at zero and
has a negative quadratic term preventing positive blow-up. The other two
are linear integrations with bounded coefficients.

For deterministic $h$, the classical controlled linear-Gaussian filter
has conditional mean $m_t=\mathbb E[X_t\mid\mathcal F_t^Y]$ and
conditional variance $P_t$ given by
\[
 dm=(am+bu)dt+hP\,dI,\qquad
 \dot P=2aP+\sigma^2-\alpha P^2,
\]
where $I_t=Y_t-\int_0^t h_sm_sds$ is an innovation Brownian motion.
We use this standard filtering identity, also the starting point of
\cite[Section 3.1]{aid}; deterministic precision makes $P$ independent
of the receiver control. Controls are observation adapted with finite
second moments and a well-posed state/observation law.

\begin{theorem}[Unique deterministic-device best response]
The receiver's unique best response, up to $dt\,d\Pr$ equality, is
\[
 u_t^*=-{b\over r_R}(K_tm_t+s_t).
\]
Its minimum cost is
\[
 \mathbb E[K_0X_0^2+2s_0X_0+d_0]
       +\int_0^T{b^2\over r_R}K_t^2P_t\,dt.
\]
\end{theorem}
\begin{proof}
Apply It\^o's formula to $K_tX_t^2+2s_tX_t+d_t$ and substitute the three
backward equations. After integrating and using the zero terminal values,
\[
 J_R=\mathbb E[K_0X_0^2+2s_0X_0+d_0]
       +r_R\mathbb E\int_0^T
          \left(u_t+{b\over r_R}(K_tX_t+s_t)\right)^2dt.
\]
Conditioning inside the square on $\mathcal F_t^Y$ replaces it by
$(u_t+b(K_tm_t+s_t)/r_R)^2+b^2K_t^2P_t/r_R^2$.
The latter term is deterministic and control independent. The asserted
linear feedback sets the first term to zero and has a unique well-posed
linear filtered-state system with finite second moments. It attains the
lower bound. Equality for any other optimum forces the same control
almost everywhere; the linear closed-loop law then gives uniqueness.
Localization and $L^2$ estimates justify the integrals in the identity.
\end{proof}

\section*{A finite-dimensional sender problem}
Put $A_t=a-b^2K_t/r_R$, $f_t=-b^2s_t/r_R$, and let
$M_t=\mathbb E m_t$, $S_t=\operatorname{Var}(m_t)$. Since $m_0$ is
nonrandom, these satisfy
\[
 \dot M=AM+f,\quad M(0)=m_0,\qquad
 \dot S=2AS+\alpha P^2,\quad S(0)=0.
\]
The conditional variance identity gives $\operatorname{Var}(X)=P+S$.
Thus the exact deterministic-schedule sender objective is
\[
 \int_0^T\left[q_S\{(M-\ell_S)^2+P+S\}
 +{r_Sb^2\over r_R^2}\{(KM+s)^2+K^2S\}+c_S\alpha\right]dt. \tag{1}
\]
The mean path $M$ does not depend on the schedule. Only two controlled
state coordinates, $P,S$, remain.

\begin{theorem}[Attainment without a relaxed-device gap]
The infimum of (1) over measurable $\alpha:[0,T]\to[0,\bar h^2]$
is attained. Taking $h=\sqrt\alpha$ implements the minimizer as an
ordinary deterministic admissible observation device.
\end{theorem}
\begin{proof}
For every schedule, positivity and the scalar differential inequality
$\dot P\le2|a|P+\sigma^2$ give a common finite upper bound for $P$.
The $S$ equation and Gronwall give a common upper bound for $S$ as well.
The derivatives of both states are uniformly bounded. From a minimizing
sequence extract uniformly convergent state paths by compactness, and a
weak-star convergent subsequence of the bounded $L^\infty$ controls.
The limiting control remains in $[0,\bar h^2]$ almost everywhere.
In each integral equation, the term $\alpha_nP_n^2$ converges weakly in
integrals: its difference from $\alpha_nP^2$ is bounded by a constant
times $\|P_n-P\|_\infty$, while the latter converges by weak-star
convergence tested against $P^2\mathbf1_{[0,t]}\in L^1$.
The limiting paths therefore solve the original equations. Every term of
(1) converges, including $c_S\int\alpha_n$. The limit attains the
infimum and is implemented by the bounded measurable square root. No
extra randomization or measure-valued control has been introduced.
\end{proof}

\section*{The unsolved extension}
For an observation-dependent device, the variance equation, when justified,
has a random precision, and the law of its variance path can depend on
receiver actions through future observations. In the completed-square
identity the term $\mathbb E\int K_t^2P_tdt$ is then not established to
be control independent. Minimizing only the visible square is therefore
not a proof of receiver optimality for arbitrary $H$. Nor can an HJB for
(1) represent this larger commitment problem without a sufficient state
for the receiver's incentives. The attained deterministic value is an
upper bound on the full sender infimum, not an equality claim. Resolving
this dependence, proving feedback-device best-response existence, and
characterizing the full sender value remain open in this attempt.

\begin{thebibliography}{9}
\bibitem{aid} R. A\"id, O. Bonesini, G. Callegaro and L. Campi (2024),
\emph{Continuous-time persuasion by filtering}, arXiv:2410.07735v1,
Introduction and Section 3.1. \url{https://arxiv.org/html/2410.07735v1}.
\end{thebibliography}
\end{document}
